%% ma: L12-B06-0058
%% lop: 12
%% chuong: 2
%% bai: 6
%% dang: 12.6.7
%% loai: TN4
%% muc: TH
%% dap_an: "D"
%% nguon: "CĐ2.1.pdf giáo viên gửi 27/10/2026 (M07, Lê Minh Khôi – Nguyễn Anh Khôi, Chương II Chủ đề 1 Vectơ trong không gian), Câu 18"
%% kien_thuc: "Bài 6 (quy tắc ba điểm, tâm hình bình hành)"
%% trang_thai: chinh-thuc
\begin{ex}
Cho hình lập phương $ABCD.A'B'C'D'$ có cạnh bằng $a$. Gọi $O$ là tâm của hình vuông $ABCD$ và $S$ là điểm thỏa mãn
\[\overrightarrow{OS}=\overrightarrow{OA}+\overrightarrow{OB}+\overrightarrow{OC}+\overrightarrow{OD}+\overrightarrow{OA'}+\overrightarrow{OB'}+\overrightarrow{OC'}+\overrightarrow{OD'}.\]
Độ dài $OS$ bằng
\choice{$OS=6a$}{$OS=2a$}{$OS=a$}{\True $OS=4a$}
\loigiai{Gọi $O'$ là tâm hình vuông $A'B'C'D'$. Ta có $\overrightarrow{OA}+\overrightarrow{OB}+\overrightarrow{OC}+\overrightarrow{OD}=\vec 0$ và $\overrightarrow{OA'}+\overrightarrow{OB'}+\overrightarrow{OC'}+\overrightarrow{OD'}=4\overrightarrow{OO'}$.

Vậy $OS=4OO'=4a$.}
\end{ex}
